\section{Internet das Coisas (IoT), IIoT e Sistemas Inteligentes Distribuídos}

\textbf{OBJETIVO GERAL:}

Desenvolver competências para projetar, desenvolver, integrar e implantar soluções baseadas em Internet das Coisas (IoT), Internet Industrial das Coisas (IIoT), Edge Computing e Computação Distribuída, aplicando Inteligência Artificial em ambientes industriais, corporativos e cidades inteligentes.

\textbf{CARGA HORÁRIA:} 120 horas

\textbf{FUNÇÃO:}

Desenvolver sistemas computacionais baseados em Inteligência Artificial, Análise Preditiva e Ecossistemas IoT, atendendo às normas de qualidade corporativas, robustez metodológica, governança de dados e arquitetura segura.

\textbf{SUBFUNÇÕES:}

\begin{itemize}
\item Projetar arquiteturas de Internet das Coisas.
\item Desenvolver soluções utilizando dispositivos IoT e IIoT.
\item Integrar sensores, gateways e plataformas inteligentes.
\item Implementar sistemas distribuídos para coleta e processamento de dados.
\item Aplicar Inteligência Artificial em ambientes IoT.
\end{itemize}

\textbf{PADRÕES DE DESEMPENHO:}

\begin{itemize}
\item Identificar arquiteturas IoT e IIoT.
\item Integrar dispositivos inteligentes utilizando protocolos industriais.
\item Desenvolver aplicações para aquisição de dados.
\item Implementar sistemas distribuídos para processamento em tempo real.
\item Aplicar Edge Computing em soluções inteligentes.
\item Desenvolver pipelines para telemetria.
\item Aplicar requisitos de segurança em ambientes IoT.
\item Monitorar dispositivos conectados.
\end{itemize}

\textbf{CAPACIDADES TÉCNICAS:}

\textbf{Arquiteturas IoT}

\begin{itemize}
\item Interpretar arquiteturas IoT e IIoT.
\item Identificar componentes de sistemas distribuídos.
\item Projetar soluções para aquisição de dados.
\item Modelar arquiteturas escaláveis.
\end{itemize}

\textbf{Dispositivos Inteligentes}

\begin{itemize}
\item Configurar sensores.
\item Integrar atuadores.
\item Configurar gateways.
\item Desenvolver aplicações embarcadas básicas.
\end{itemize}

\textbf{Comunicação}

\begin{itemize}
\item Utilizar protocolos MQTT.
\item Utilizar OPC UA.
\item Utilizar CoAP.
\item Integrar dispositivos em redes distribuídas.
\end{itemize}

\textbf{Processamento Distribuído}

\begin{itemize}
\item Desenvolver pipelines de telemetria.
\item Processar dados em tempo real.
\item Utilizar bancos de séries temporais.
\item Integrar Edge e Cloud Computing.
\end{itemize}

\textbf{Segurança}

\begin{itemize}
\item Aplicar mecanismos de autenticação.
\item Proteger dispositivos conectados.
\item Garantir integridade dos dados.
\item Aplicar princípios de segurança industrial.
\end{itemize}

\textbf{CONHECIMENTOS:}

\textbf{Fundamentos de IoT}

\begin{itemize}
\item Internet das Coisas
\item Internet Industrial das Coisas (IIoT)
\item Arquiteturas IoT
\item Arquiteturas IIoT
\item Ecossistemas Inteligentes
\end{itemize}

\textbf{Dispositivos}

\begin{itemize}
\item Sensores
\item Atuadores
\item Controladores
\item Gateways
\item Dispositivos Embarcados
\end{itemize}

\textbf{Protocolos}

\begin{itemize}
\item MQTT
\item OPC UA
\item CoAP
\item HTTP
\item WebSockets
\end{itemize}

\textbf{Modelagem de Dados}

\begin{itemize}
\item JSON
\item Avro
\item Protobuf
\item Serialização
\item Estruturas de Mensagens
\end{itemize}

\textbf{Computação Distribuída}

\begin{itemize}
\item Edge Computing
\item Fog Computing
\item Cloud Computing
\item Sistemas Distribuídos
\item Alta Disponibilidade
\end{itemize}

\textbf{Streaming de Dados}

\begin{itemize}
\item Apache Kafka
\item Apache Spark
\item Apache Flink
\item Processamento em Tempo Real
\item Stream Processing
\end{itemize}

\textbf{Bancos de Dados}

\begin{itemize}
\item InfluxDB
\item TimescaleDB
\item Banco de Séries Temporais
\item Retenção de Dados
\end{itemize}

\textbf{Arquiteturas Industriais}

\begin{itemize}
\item IIRA
\item RAMI 4.0
\item ISA-95 (conceitos)
\item Indústria 4.0
\item Digital Twin (conceitos)
\end{itemize}

\textbf{Integração com IA}

\begin{itemize}
\item Machine Learning em IoT
\item Edge AI
\item IA para Manutenção Preditiva
\item IA para Monitoramento Inteligente
\item Inferência em Borda
\end{itemize}

\textbf{Segurança e Governança}

\begin{itemize}
\item LGPD
\item Segurança em IoT
\item Criptografia
\item Controle de Acesso
\item Gestão de Dispositivos
\end{itemize}

\textbf{CAPACIDADES SOCIOEMOCIONAIS:}

\begin{itemize}
\item Demonstrar responsabilidade na integração de dispositivos conectados.
\item Trabalhar colaborativamente em projetos multidisciplinares.
\item Avaliar criticamente arquiteturas distribuídas.
\item Demonstrar organização durante a implantação de soluções IoT.
\item Produzir documentação técnica de arquiteturas desenvolvidas.
\item Compartilhar conhecimento sobre tecnologias emergentes.
\item Atuar observando princípios de segurança da informação.
\item Aplicar boas práticas de governança em ecossistemas inteligentes.
\end{itemize}

\textbf{AMBIENTES PEDAGÓGICOS:}

\begin{itemize}
\item Laboratório de IoT.
\item Laboratório de Redes.
\item Laboratório de Inteligência Artificial.
\item Ambiente Virtual de Aprendizagem.
\item Ambiente Cloud para integração de dispositivos.
\end{itemize}

\textbf{EQUIPAMENTOS E RECURSOS:}

\textbf{Hardware}

\begin{itemize}
\item Computadores com acesso à Internet.
\item Kits IoT (ESP32, Raspberry Pi ou equivalentes).
\item Sensores e atuadores.
\item Gateways IoT.
\item Projetor multimídia.
\end{itemize}

\textbf{Software}

\begin{itemize}
\item Python
\item Node-RED
\item Mosquitto MQTT
\item InfluxDB
\item TimescaleDB
\item Docker
\item Kubernetes
\item VS Code
\item Grafana
\item Git
\end{itemize}
